\par
Wir werden uns von nun ab stets der Thetafunktionen \(\vartheta_{gh}(\psi|x)\) bedienen,
jedoch der kürzern Schreibweise wegen für ihre Nullwerte die \textsc{Gaussschen}
Zeichen \(p,q,r\), also
\[
\vartheta_{00}(0|x)=p(x),\quad \vartheta_{01}(0|x)=q(x),\quad \vartheta_{10}(0|x)=r(x)
\]
\noindent setzen. Es ist dann
\begingroup\makeatletter\tagsleft@true\let\veqno\leqno\makeatother
\begin{equation*}
\left\{\begin{aligned}
W\varphi&=\sqrt{M(1,\cos v)}\,\vartheta_{00}(\psi|x),\\[1ex]
T\varphi&=\frac1{\sqrt\mu}\sqrt{M(1,\cos v)}\,\vartheta_{11}(\psi|x),\\[1ex]
\operatorname{Sl}\varphi&=\frac1{\sqrt\mu}\frac{\vartheta_{11}(\psi|x)}{\vartheta_{00}(\psi|x)}\mathpunct{\vcenter{\hbox{.}}}\\[1ex]
T(\tfrac12\varpi-\varphi)&=\frac1{\sqrt\mu}\sqrt{M(1,\cos v)}\,\vartheta_{10}(\psi|x),\\[1ex]
W(\tfrac12\varpi-\varphi)&=\sqrt{M(1,\cos v)}\,\vartheta_{01}(\psi|x),
\end{aligned}\right.\tag{22}
\end{equation*}
\endgroup
\clearpage\SourceLeaf{162}\ArithmeticHead{77}{R}{IVb. ALLGEMEINE ELLIPTISCHE FUNKTIONEN.}
\indent \textsc{Gauss} macht nun den Ansatz (Werke X, 1, S. 195, Art. 3)
\begingroup\makeatletter\tagsleft@true\let\veqno\leqno\makeatother
\begin{equation*}
\left\{\begin{aligned}
\vartheta_{11}(\psi|x)^2&=A\vartheta_{00}(2\psi|x^2)+B\vartheta_{10}(2\psi|x^2),\\
\vartheta_{00}(\psi|x)^2&=C\vartheta_{00}(2\psi|x^2)+D\vartheta_{10}(2\psi|x^2).
\end{aligned}\right.\tag{22\,a}
\end{equation*}
\endgroup
\noindent Er hat also erkannt, dass die Quadrate der \(\vartheta_{gh}(\psi|x)\) und die \(\vartheta_{gh}(2\psi|x^2)\) einem
gemeinsamen Funktionstypus angehören — man bezeichnet diesen gegenwärtig
als Thetafunktionen zweiter Ordnung —, und dass zwischen je drei zusammengehörigen dieser Funktionen eine homogene lineare Relation mit von \(\psi\) unabhängigen Koeffizienten besteht. Er bringt also hier, für den Fall der
zweiten Ordnung, das Prinzip zur Anwendung, das man gewöhnlich auf \textsc{Hermite} zurückführt\footnote{Siehe \textsc{Hermites} Brief an \textsc{Jacobi} vom August 1844, \textsc{Jacobis} Werke II, S. 96, \textsc{Hermite} Oeuvres I, S. 24, 25. Noch deutlicher tritt dieses Prinzip in einer im Handbuch 16, Bb, aufgezeichneten, aus 1808 stammenden Abhandlung auf (Werke X, 1, S. 287, besonders Art. 5, 11, 12), vergl. auch die Abhandlung aus Handbuch 18, Bd, die Werke III, S. 416 abgedruckt ist und aus 1809 stammt, besonders S. 457. Dass die letztere Rechnung auf das \textsc{Hermitesche} Prinzip gegründet ist, hat P. \textsc{Günther} in seinem Habilitationsvortrag (Göttinger Nachrichten 1894, S. 102) hervorgehoben.}. — Die Koeffizienten \(A,B,C,D\) bestimmt \textsc{Gauss} — wohl
indem er für \(\psi\) spezielle Werte einsetzt — in ihrer Abhängigkeit von \(\mu=\operatorname{tg}v\).
Den Zusammenhang, der zwischen dieser Grösse und den Thetanullwerten
\(p,q,r\) besteht, gibt er in der Formel (Werke X, 1, S. 194, Gl. (3))
\begingroup\makeatletter\tagsleft@true\let\veqno\leqno\makeatother
\begin{equation*}
T(90^{0})=\sqrt{\cos v},\tag{23}
\end{equation*}
\endgroup
\noindent die, da hier \(90^{0}\) (\textsc{Gauss} denkt sozusagen lemniskatisch!) im Bogenmass gleich
\(\tfrac12\varpi\) zu nehmen ist, mit
\[
\frac1{\sqrt{\operatorname{tg}v}}\sqrt{M(1,\cos v)}\,r(x)=\sqrt{\cos v},
\]
\noindent oder ins Quadrat erhoben mit
\begingroup\makeatletter\tagsleft@true\let\veqno\leqno\makeatother
\begin{equation*}
\frac{\sin v}{M(1,\cos v)}=r(x)^2\tag{23\,a}
\end{equation*}
\endgroup
\noindent übereinstimmt, also, wenn man beachtet, dass unser
\[
x=e^{-\pi\frac{\varpi'}{\varpi}}=e^{-\pi\frac{M(1,\cos v)}{M(1,\sin v)}}
\]
\noindent mit dem im \textit{Leiste} benutzten \(z\) (siehe den Abschnitt IVa, S. 65) durch
\(z=x^{\frac12}\) verbunden ist, nichts anderes ist, als die Leistegleichung (5). Wir
wollen bei der Wiedergabe der \textsc{Gaussschen} Formeln gleich die von \(\psi\) unabhängigen Faktoren durch die \(p,q,r\) ausdrücken und schreiben zunächst die\par
\clearpage\SourceLeaf{163}\ArithmeticHead{78}{L}{L. SCHLESINGER, ÜBER GAUSS’ FUNKTIONENTHEORETISCHE ARBEITEN.}
\noindent beiden Gleichungen (Werke X, 1, S. 196, (14), (15))
\begingroup\makeatletter\tagsleft@true\let\veqno\leqno\makeatother
\begin{align*}
2p(x^2)\vartheta_{00}(2\psi|x^2)&=\vartheta_{00}(\psi|x)^2+\vartheta_{01}(\psi|x)^2,\tag{24}\\
2p(x^2)\vartheta_{10}(2\psi|x^2)&=\vartheta_{10}(\psi|x)^2-\vartheta_{11}(\psi|x)^2\tag{25}
\end{align*}
\endgroup
\noindent hin, die \textsc{Gauss} aus den (22\,a) entsprechenden ableitet, indem er die durch das
gedachte \textsc{Hermitesche} Prinzip gewonnenen Ausdrücke der \(\vartheta_{10}^{2},\vartheta_{01}^{2}\) durch die
\(\vartheta_{00}^{2},\vartheta_{11}^{2}\) (a. a. O. Gln. (16))
\begingroup\makeatletter\tagsleft@true\let\veqno\leqno\makeatother
\begin{equation*}
\left\{\begin{aligned}
p(x)^2\vartheta_{10}(\psi|x)^2&=r^2\vartheta_{00}(\psi|x)^2-q^2\vartheta_{11}(\psi|x)^2,\\
p(x)^2\vartheta_{01}(\psi|x)^2&=q^2\vartheta_{00}(\psi|x)^2+r^2\vartheta_{11}(\psi|x)^2
\end{aligned}\right.\tag{26}
\end{equation*}
\endgroup
\noindent heranzieht. — In gleicher Weise ergibt sich die Gleichung (a. a. O. Gl. (17))
\begingroup\makeatletter\tagsleft@true\let\veqno\leqno\makeatother
\begin{equation*}
2\vartheta_{01}(\psi|x)\vartheta_{11}(\psi|x)=r(\sqrt x)\vartheta_{11}(\psi|\sqrt x),\tag{27}
\end{equation*}
\endgroup
\noindent die zu der sogenannten \textsc{Gaussschen} Transformation gehört, während die Gleichungen (24), (25) die \textsc{Landensche} Transformation darstellen, auf die wir alsbald ausführlich zurückzukommen haben. Hier geben wir noch die übrigen
zu diesen Transformationen gehörigen Gleichungen an, die zwar an dieser
Stelle bei \textsc{Gauss} fehlen, aber in ganz analoger Weise abgeleitet werden
können\footnote{Siehe bei \textsc{Gauss}, Werke III, S. 471, 472, bei \textsc{Weber} a. a. O.\textsuperscript{102)}, § 18, 27.}:
\begingroup\makeatletter\tagsleft@true\let\veqno\leqno\makeatother
\begin{align*}
q(x^2)\vartheta_{01}(2\psi|x^2)&=\vartheta_{00}(\psi|x)\vartheta_{01}(\psi|x)\tag{25\,a}\\
\vartheta_{00}(\psi|x)^2+\vartheta_{10}(\psi|x)^2&=p(\sqrt x)\vartheta_{00}(\psi|\sqrt x)\tag{27\,a}\\
2\vartheta_{00}(\psi|x)\vartheta_{10}(\psi|x)&=r(\sqrt x)\vartheta_{10}(\psi|\sqrt x).\tag{27\,b}
\end{align*}
\endgroup

Diese Formeln liefern nun zunächst den in der Tagebuchnotiz 105. angezeigten Zusammenhang mit der Theorie des agM. Indem man nämlich
\(\psi=0\) setzt, erhält man aus (24), (25), (25\,a), (27\,a) die (bei \textsc{Gauss} in einer
aus 1825 stammenden Abhandlung \textit{Hundert Theoreme etc.} Werke III, S. 466
aufgezeichneten) Formeln:
\begingroup\makeatletter\tagsleft@true\let\veqno\leqno\makeatother
\begin{equation*}
\left\{\begin{aligned}
2p(x^2)^2&=p(x)^2+q(x)^2,\\
2p(x^2)r(x^2)&=r(x)^2,\\
q(x^2)^2&=p(x)q(x),\\
p(x^2)^2+r(x^2)^2&=p(x)^2,\\
p(x^2)^2-r(x^2)^2&=q(x)^2.
\end{aligned}\right.\tag{28}
\end{equation*}
\endgroup
\clearpage\SourceLeaf{164}\ArithmeticHead{79}{R}{IVb. ALLGEMEINE ELLIPTISCHE FUNKTIONEN.}
\noindent Setzen wir also
\begingroup\makeatletter\tagsleft@true\let\veqno\leqno\makeatother
\begin{equation*}
\left\{\begin{aligned}
a&=p(x)^2,& b&=q(x)^2,& c&=r(x)^2,\\
a_1&=p(x^2)^2,& b_1&=q(x^2)^2,& c_1&=r(x^2)^2,
\end{aligned}\right.\tag{29}
\end{equation*}
\endgroup
\noindent so bestehen (vergl. S. 65) zwischen diesen Grössen die Relationen des agM:
\begingroup\makeatletter\tagsleft@true\let\veqno\leqno\makeatother
\begin{equation*}
\left\{\begin{aligned}
c^2&=a^2-b^2,& c_1^2&=a_1^2-b_1^2,\\
a_1&=\tfrac12(a+b),& b_1&=\sqrt{ab},\\
a&=a_1+c_1,& b&=a_1-c_1.
\end{aligned}\right.\tag{28\,a}
\end{equation*}
\endgroup
\noindent Wir erhalten daher, wenn wir \(x\) der Reihe nach zur 4., 8., \(\ldots\) Potenz erheben und das agM. bilden\footnote{Siehe bei \textsc{Gauss}, Werke III, S. 467.}
\begingroup\makeatletter\tagsleft@true\let\veqno\leqno\makeatother
\begin{equation*}
M(a,b)=M(p(x)^2,\ q(x)^2)=\lim_{n\to\infty}p(x^{2^n})^2=1,\tag{30}
\end{equation*}
\endgroup
\noindent wo wir die bisher stillschweigend gemachte Voraussetzung, dass \(x\) dem absoluten Betrage nach kleiner ist als 1, ausgenutzt haben.

In der Tat lag der angedeuteten Aufstellung des obigen Formelsystems
die Grösse \(x\) zugrunde, die nur der Beschränkung zu unterwerfen ist, dass
ihr absoluter Betrag kleiner als 1 sei, damit die Konvergenz der Reihen
\(\vartheta_{gh}(\psi|x)\) gesichert erscheint. — Bei \textsc{Gauss} tritt dadurch, dass er von einer
Grösse \(\mu=\operatorname{tg}v\) ausgeht, noch eine Schwierigkeit hinzu, auf die wir jetzt,
ehe wir weiter gehn, hinweisen wollen. — \textsc{Gauss} bildet mit dem gegebenen
\(\operatorname{tg}v\) zunächst die agM. \(M(1,\cos v)\) und \(M(1,\sin v)\) und dann weiter die durch
die Gleichungen (18) gegebenen Grössen
\begingroup\makeatletter\tagsleft@true\let\veqno\leqno\makeatother
\begin{equation*}
\varpi=\frac{\pi\cos v}{M(1,\cos v)}\mathpunct{\vcenter{\hbox{,}}}\quad
\varpi'=\frac{\pi\cos v}{M(1,\sin v)}\mathpunct{\vcenter{\hbox{.}}}\tag{18\,a}
\end{equation*}
\endgroup
\noindent Solange diese positiv sind, ist das mit denselben gebildete \(x\) positiv und kleiner
als 1, also die Konvergenz der Reihen \(\vartheta_{gh}(\psi|x)\) gesichert. Die oben skizzierte
Theorie des agM. liefert dann neben der Gleichung (23\,a) noch die Gleichungen
\begingroup\makeatletter\tagsleft@true\let\veqno\leqno\makeatother
\begin{equation*}
1=M(1,\cos v)p(x)^2,\quad\cos v=M(1,\cos v)q(x)^2,\tag{23\,b}
\end{equation*}
\endgroup
\noindent die in Verbindung mit (23\,a) auch die Bestimmung
\begingroup\makeatletter\tagsleft@true\let\veqno\leqno\makeatother
\begin{equation*}
\mu=\operatorname{tg}v=\frac{r(x)^2}{q(x)^2}\tag{23\,c}
\end{equation*}
\endgroup
\clearpage\SourceLeaf{165}\ArithmeticHead{80}{L}{L. SCHLESINGER, ÜBER GAUSS’ FUNKTIONENTHEORETISCHE ARBEITEN.}
\noindent ergeben\footnote{Die Gleichungen (16) oben S. 69 lauten, wenn wir \(\nu=\sin v\) nehmen,
\[
\begin{gathered}
\sin v=M(1,\cos v)\,.\,2p(x^2)r(x^2),\\
1=M(1,\cos v)(p(x^2)^2+r(x^2)^2),
\end{gathered}
\]
was aber gemäss den Formeln (28) mit den Gleichungen (23\,a, b) übereinstimmt.}.
Für ein beliebiges reelles oder komplexes \(\mu\) entsteht hier zunächst die Aufgabe, zu beweisen, dass der reelle Teil von \(M(1,\cos v)/M(1,\sin v)\)
stets wesentlich positiv ist, vorausgesetzt, dass unter den unendlich vielen
Werten der agM. jeweils zwei »zusammengehörige« zur Bildung dieses Quotienten benutzt werden. Wie \textsc{Gauss} diese Frage bewältigt hat, werden wir
nachher auseinandersetzen; vorerst wollen wir, um von der hierdurch bezeichneten Schwierigkeit unbelastet, den durch die Tagebuchnotiz 106. gewiesenen
Weg weiter verfolgen zu können, die Funktion \(\operatorname{Sl}\varphi\) in der Weise definiert
denken, dass wir gemäss den Gleichungen (23\,a)—(23\,c), die ja überhaupt den
Übergang von den bei \textsc{Gauss} auftretenden Ausdrücken der Koeffizienten in
den Gleichungen (24)—(27), zu den von uns benutzten vermitteln, den bei
Sl auftretenden Faktor \(1/\sqrt\mu\) (siehe die Gl. (22)) durch \(q(x)/r(x)\) ersetzen und
den in \(\varphi=\psi\varpi\) auftretenden Faktor \(\varpi\) vorerst unbestimmt lassen. — Wir
gehen also von der Definition
\begingroup\makeatletter\tagsleft@true\let\veqno\leqno\makeatother
\begin{equation*}
\operatorname{Sl}\psi\varpi=\frac{q(x)}{r(x)}\frac{\vartheta_{11}(\psi|x)}{\vartheta_{00}(\psi|x)}\tag{31}
\end{equation*}
\endgroup
\noindent aus und lesen nunmehr in der Tagebuchnotiz 108. von Ende Mai den ersten
Satz: »Es ist gelungen Zähler und Nenner des Sinus lemniscaticus (in allgemeinster Fassung) auf Integralgrössen zu reduzieren«. — Die Bedeutung dieses
Satzes kann nichts anderes sein, als die Verifikation, dass Sl tatsächlich einer
Differentialgleichung von der Form
\begingroup\makeatletter\tagsleft@true\let\veqno\leqno\makeatother
\begin{equation*}
\frac{du}{d\varphi}=\sqrt{1+\mu^2\sin^2u}\tag{32}
\end{equation*}
\endgroup
\noindent Genüge leistet, wenn man
\begingroup\makeatletter\tagsleft@true\let\veqno\leqno\makeatother
\begin{equation*}
\sin u=\operatorname{Sl}\psi\varpi\tag{31\,a}
\end{equation*}
\endgroup
\noindent setzt. — In der Scheda Ac findet sich eine solche Verifikation nicht, auch
in keiner aus derselben Zeit stammenden anderen Aufzeichnung, wohl aber
hat sie \textsc{Gauss} in einer 1825 verfassten Notiz des Handbuchs 16, Bb, durchgeführt (Werke X, 1, S. 308—310), und diese Rechnung ordnet sich schon
darum ganz ungezwungen hier ein, weil sie auch wesentlich auf dem Prinzip
beruht, das der Herleitung der Formeln (24) bis (27) zugrunde liegt. Es ist\par
\clearpage\SourceLeaf{166}\ArithmeticHead{81}{R}{IVb. ALLGEMEINE ELLIPTISCHE FUNKTIONEN.}
\HistoricalFoot{X\textsubscript{2} Abh. 2.}{11}
\noindent nämlich der bei der Differentiation von Sl auftretende Zähler \(\vartheta_{11}'\vartheta_{00}-\vartheta_{00}'\vartheta_{11}\)
(wo der Akzent die Differentiation nach \(\psi\) bedeutet und, wie auch im folgenden, bei den \(\vartheta_{gh}\) die Argumente \(\psi,x\) unterdrückt werden) eine Thetafunktion zweiter Ordnung, die sich von dem Produkt \(\vartheta_{10}\vartheta_{01}\) nur durch einen von
\(\psi\) unabhängigen Faktor unterscheidet. Über die Bestimmung dieses Faktors
wird nachher noch einiges zu bemerken sein; es ergibt sich\footnote{Siehe \textsc{Weber}, a. a. O.\textsuperscript{102)}, § 20, Gl. (7).}
\begingroup\makeatletter\tagsleft@true\let\veqno\leqno\makeatother
\begin{equation*}
\vartheta_{11}'\vartheta_{00}-\vartheta_{00}'\vartheta_{11}=\pi p(x)^2\vartheta_{10}\vartheta_{01}.\tag{33}
\end{equation*}
\endgroup
\noindent Erhebt man diese Gleichung ins Quadrat und drückt die rechter Hand auftretenden Quadrate der \(\vartheta_{10},\vartheta_{01}\) mittels der von \textsc{Gauss} in der Scheda Ac aufgezeichneten Gleichungen (26) aus, so findet man die in der Notiz von 1825
(Werke X, 1, S. 310) von \textsc{Gauss} auf etwas andere Weise\footnote{\textsc{Gauss} betrachtet a. a. O. (vergl. \textsc{Geppert}, a. a. O.\textsuperscript{83)}, § 24, S. 95) anstelle des Ausdrucks (33) die die zweiten Derivierten \(\vartheta_{gh}''\) nach \(\psi\) enthaltenden Ausdrücke
\[
\vartheta_{gh}^{\prime 2}-\vartheta_{gh}\vartheta_{gh}''
\]
für \(g,h=0,0\) und \(1,1\). Diese sind ebenfalls Thetafunktionen zweiter Ordnung und lassen sich linear durch die Quadrate der \(\vartheta_{00},\vartheta_{11}\) darstellen (vergl. \textsc{Weber}, a. a. O.\textsuperscript{102)}, § 20). Dass sich dieser Vorgang von dem des Textes nicht wesentlich unterscheidet, erhellt daraus, dass die \(\vartheta_{gh}\) die partielle Differentialgleichung
\[
\frac{\partial^2\vartheta}{\partial\psi^2}+4\pi^2x\frac{\partial\vartheta}{\partial x}=0
\]
befriedigen (\textsc{Weber}, a. a. O., § 17, (4)). Ein Unterschied ergibt sich nur bei der Bestimmung der von \(\psi\) unabhängigen Faktoren, und dies dürfte auch der Grund sein, weshalb \textsc{Gauss} an dieser Stelle nicht direkt mit (33) operiert hat.} hergeleitete Gleichung:
\begingroup\makeatletter\tagsleft@true\let\veqno\leqno\makeatother
\begin{equation*}
(\vartheta_{11}'\vartheta_{00}-\vartheta_{00}'\vartheta_{11})^2
=\pi^2(r^2\vartheta_{00}^2-q^2\vartheta_{11}^2)(q^2\vartheta_{00}^2+r^2\vartheta_{11}^2).\tag{34}
\end{equation*}
\endgroup
\noindent Nun differenziert \textsc{Gauss} die Gleichung (31\,a) und erhält mit Rücksicht auf
\[
\cos u=\frac1{r\vartheta_{00}}(r^2\vartheta_{00}^2-q^2\vartheta_{11}^2)^{\frac12}
\]
\noindent die Differentialgleichung
\begingroup\makeatletter\tagsleft@true\let\veqno\leqno\makeatother
\begin{equation*}
\frac{du}{d\varphi}=\frac{\pi q^2}{\varpi}\sqrt{1+\frac{r^4}{q^4}\sin^2u}\text{\textsuperscript{109)}}.\tag{35}
\end{equation*}
\endgroup
\setcounter{footnote}{109}
\footnotetext{\textsc{Gauss} hat nur (Werke X, 1, S. 310) statt des Faktors \(q(x)/r(x)\), den wir in (31) geschrieben haben (vergl. \textsc{Geppert} a. a. O.\textsuperscript{83)}, S. 146), den Faktor
\[
\sqrt{\frac{p(x^2)^2-r(x^2)^2}{2p(x^2)r(x^2)}},
\]
der aber vermöge der Gleichungen (28) mit dem unsern übereinstimmt; vergl. auch die Fussnote 106).}
\clearpage\SourceLeaf{167}\ArithmeticHead{82}{L}{L. SCHLESINGER, ÜBER GAUSS’ FUNKTIONENTHEORETISCHE ARBEITEN.}
\noindent Setzt man hierin
\begingroup\makeatletter\tagsleft@true\let\veqno\leqno\makeatother
\begin{equation*}
\varpi=\pi q^2,\quad\frac{r^4}{q^4}=\mu^2\tag{36}
\end{equation*}
\endgroup
\noindent so erhält (35) die Form (32), und damit ist das Ziel erreicht.

Wir wenden uns nun vorerst zu der Frage, nach der Bestimmung der von
\(\psi\) unabhängigen Faktoren, die in den durch das \textsc{Hermitesche} Prinzip gegebenen
Ansätzen auftreten. — Diese Bestimmung wird sehr wesentlich erleichtert,
wenn man die Darstellung der \(\vartheta_{gh}\) und ihrer Nullwerte \(p,q,r\) in der Form
\textls[100]{unendlicher Produkte} heranzieht. Hierauf bezügliche Rechnungen finden
wir in der Scheda Ac von Seite 40 an (Werke X 1, S. 201ff., Art. 9—14);
sie gipfeln in der berühmten Identität zwischen den Darstellungen der Thetafunktion als unendliche Reihe und als unendliches Produkt, die wir oben
(Abschnitt II, Gl. (3), S. 16) aus den \textit{Fundamenta} \textsc{Jacobis} reproduziert haben,
und die sich bei \textsc{Gauss} hier im Art. 12 (a. a. O., S. 204) findet. — \textsc{Gauss} ist
später noch wiederholt auf den Zusammenhang zwischen Produkt- und Reihendarstellung der \(\vartheta_{gh}\) zurückgekommen, namentlich hat er jene Identität wiederholt abgeleitet\footnote{Siehe Werke X, 1, S. 287, besonders Art. 4, S. 292, 293, ferner \textit{Hundert Theoreme}, Werke III, S. 461, besonders Art. 4, S. 464, das »vierte Theorem«.}. Wahrscheinlich hat er sie auch 1800 schon in ähnlicher
Weise bewiesen, wie später in den \textit{Hundert Theoremen}, nämlich durch Grenzübergang aus einer Identität zwischen einer endlichen Summe und einem endlichen Produkt (Werke III, S. 462, das »erste Theorem«), also ebenso wie
später \textsc{Cauchy}. — In der Scheda folgt dann (Werke X, 1, S. 205, Art. 13)
die Aufstellung der Nullstellen für die Funktionen \(\vartheta_{11}\) und \(\vartheta_{00}\), sowie (a. a. O.,
S. 206, Art. 14) ihre Produktdarstellung in der definitiven Form.

Man wird diese Entwicklungen mit dem zweiten Satze der Tagebuchnotiz 108. in Verbindung zu setzen haben: »auch die aus genuinen Prinzipien
geschöpfte Entwicklung aller nur denkbaren lemniskatischen Funktionen in
unendliche Reihen; eine überaus schöne Erfindung, die wohl keiner der vorhergehenden nachsteht«. In der Tat ist \textsc{Gauss} auf Grund jener Identität imstande, die aus der Kenntnis der Nullstellen einer beliebigen Thetafunktion
sich ergebende Produktentwicklung in eine Reihenentwicklung umzuwandeln.
Wie elegant sich mit Hilfe der Produktdarstellung die Faktorenbestimmung
in den Gln. (24)—(27) gestaltet, zeigt die 1827 ins Handbuch 16, Bb, eingetragene Rechnung (Werke III, S. 470).
